\section{Quantitative Study}
\label{sec:experiments}

\begin{figure}[t]
  \centering
  \placeholderfig{32mm}{figures/qualitative.png --- three or four examples:
  original image with detected box, unmodified prediction, intervened
  prediction, and the resulting displacement in km}
  \caption{Qualitative examples. \TODO{describe each case}}
  \label{fig:qualitative}
  \vspace{-3mm}
\end{figure}

\subsection{Experimental Setup}

We evaluate on Img2GPS3k~\cite{vo2017revisiting}, comprising $2{,}997$
images with ground-truth coordinates. Grounding DINO produces the regions
from the eight-concept vocabulary of Section~\ref{sec:system}; images for
which no detected box survives the mapping into CLIP's centre crop are
excluded, and their number reported, so that the evaluated subset is
unambiguous (\TODO{$N$ evaluated, $N$ excluded}). We report mean and median
top-$1$ geodesic error and accuracy within the conventional $1$, $25$,
$200$, $750$ and $2{,}500$~km thresholds. The location encoder and gallery
embeddings are computed once and reused, so the rows of
Table~\ref{tab:main} differ solely in the attention of the vision tower.

\subsection{Grouped Attention Search}

The $24$ layers are partitioned into three contiguous groups --- early
($0$--$7$), middle ($8$--$15$) and late ($16$--$23$) --- each governed by a
single parameter $\delta_i$, with in-region bias $a_i = \delta_i$ and
out-of-region bias $b_i = -\delta_i$ for every layer of group $i$. Positive
$\delta$ favours the detected regions, negative $\delta$ favours their
complement, and $\delta_1 = \delta_2 = \delta_3 = 0$ recovers the unmodified
model exactly. Candidates are sampled uniformly from
$\delta_i \in [-3, 3]$ over \TODO{$T$} trials with seed $42$, minimising the
mean top-$1$ distance.

\begin{table}[t]
\centering
\caption{Geolocation accuracy on Img2GPS3k under attention intervention. The
first row is the unmodified GeoCLIP model. \TODO{populate from
\texttt{random\_search\_results.csv} and \texttt{best\_config.json}}}
\label{tab:main}
\setlength{\tabcolsep}{4pt}
\begin{tabular}{L{3.7cm}|P{1.2cm}P{1.2cm}|P{0.9cm}P{0.9cm}P{0.9cm}P{0.9cm}P{1cm}}
\toprule
\multirow{2}{*}{\textbf{Configuration}} &
\multicolumn{2}{c|}{\textbf{Error (km)}} &
\multicolumn{5}{c}{\textbf{Accuracy (\%) within}} \\
 & Mean & Median & 1 km & 25 km & 200 km & 750 km & 2500 km \\
\midrule
GeoCLIP (baseline) & -- & -- & -- & -- & -- & -- & -- \\
\midrule
\quad + suppress regions & -- & -- & -- & -- & -- & -- & -- \\
\quad + amplify regions & -- & -- & -- & -- & -- & -- & -- \\
\quad + searched $\delta$ & -- & -- & -- & -- & -- & -- & -- \\
\midrule
\quad + random regions (control) & -- & -- & -- & -- & -- & -- & -- \\
\bottomrule
\end{tabular}
\end{table}

\subsection{Results}

\TODO{to be written once the search has completed}

% ---------------------------------------------------------------------
% Author note (not rendered). All three outcomes support a demonstration
% paper; select the framing that matches the numbers obtained.
%
%  (a) Intervention improves accuracy. Report the gain and observe that the
%      mechanism which explains the model also steers it: the detected cues
%      carry signal that the unmodified attention under-weights.
%  (b) Intervention degrades accuracy substantially. This is the stronger
%      explainability result: the detected regions are load-bearing, and the
%      magnitude of the degradation quantifies the extent. Contrast with the
%      random-region control to establish that the effect does not follow
%      merely from perturbing attention at all.
%  (c) Little changes. Report this plainly: it bounds the proportion of
%      GeoCLIP's behaviour reachable through attention reweighting alone and
%      motivates representation-level intervention.
% ---------------------------------------------------------------------

The random-region control is material to the interpretation in every case. A
region-conditioned intervention must be compared against an area-matched
random region; otherwise an apparent effect may reflect the proportion of
the image perturbed rather than its content.
\TODO{this control is the most probable reviewer objection}
